\begin{table}[H]
    \centering
    \caption{Hardware and system prerequisites for PentaForge deployment}
    \label{tab:hardware_prerequisites}
    \renewcommand{\arraystretch}{1.35}
    \begin{tabular}{|p{0.22\textwidth}|p{0.28\textwidth}|p{0.40\textwidth}|}
        \hline
        \textbf{Component} & \textbf{Minimum} & \textbf{Recommended} \\
        \hline
        Operating system & Linux host with Docker support, preferably Ubuntu 22.04
        LTS or newer. & Ubuntu 22.04/24.04 LTS with Docker Engine and Docker Compose
        for stable backend, sandbox, Redis, Qdrant, and frontend execution. \\
        \hline
        CPU & 4 CPU cores. & 8 CPU cores or more for smoother concurrent execution
        of the FastAPI backend, sandbox tools, Qdrant vector search, and desktop UI. \\
        \hline
        RAM & 8 GB. & 16 GB or more, especially when running the embedding model,
        backend, sandbox, Redis, Qdrant, and frontend at the same time. \\
        \hline
        Storage & 50 GB SSD. & 100--200 GB SSD for Docker images, Python
        dependencies, model cache, scan artifacts, logs, reports, and sandbox
        workspaces. \\
        \hline
        GPU & Not required. & Optional only for future local LLM or accelerated
        embedding workflows. The current Docker deployment uses CPU-compatible
        PyTorch and remote or configured LLM providers. \\
        \hline
        Network & Internet connection for initial setup and target reachability. &
        Stable internet access for Docker image pulls, Python package installation,
        remote LLM calls, vulnerability intelligence updates, and live target access. \\
        \hline
    \end{tabular}
\end{table}
